\chapter{Anexos del Subdocumento 7}

Estos anexos detallan los listados que el Subdocumento 7 resume. La EDT, el diccionario y la carta Gantt están en el Formulario T-14; el método de estimación y programación, la ruta crítica y los frentes de trabajo, en el Formulario T-15; y la implantación, en el Formulario T-18.

\section{Anexo 7.A — Meses contractuales y ventanas de congelamiento según la fecha de inicio}

El Art. 17° de las Bases Administrativas fija los meses del contrato, no sus fechas. El mes 1 es el primer mes completo de ejecución (Bases Administrativas, Art. 10°, numeral 2). Si el contrato se inicia en el mes calendario m, el mes contractual n cae en el mes calendario \(m + (n - 1)\), contado en aritmética de doce meses.

El caso fija las ventanas del año: congelamiento total del 1 al 25 de septiembre y durante todo diciembre; congelamiento en los tres primeros días hábiles de cada mes (Caso 02, RT-10.05); y ningún paso a producción en septiembre ni en diciembre (sección 13.3, condición 2). El Capítulo 3, sección 3.1.2, deriva de esa regla los ocho meses de inicio que no ponen los meses 16 ni 21 en septiembre o diciembre (supuesto S-17). La Tabla \ref{tab:7A1} aplica la misma fórmula a los meses de marcha blanca; los meses en negrita caen en una ventana de congelamiento.

\parte{tablas/tab_7A1}

La tabla lleva a tres conclusiones. Primero, ningún inicio admisible deja las dos marchas blancas completas fuera de septiembre y diciembre: siempre hay al menos un mes de marcha blanca dentro de una ventana de congelamiento. Segundo, solo los inicios de diciembre de 2026, febrero de 2027 y marzo de 2027 ponen en producción la Etapa 2 antes de enero de 2029, fecha desde la que rigen las condiciones de la principal cadena de supermercados (Caso 02, sección 13.2; Capítulo 3, sección 3.1.2).

Tercero, entre esos tres inicios, febrero de 2027 es el único en que el congelamiento afecta solo el último mes de una marcha blanca (el mes 20, de la Etapa 2) y no un inicio de marcha blanca. Con diciembre de 2026, el H6 caería en diciembre; con marzo de 2027, el H11 caería en septiembre. Con febrero de 2027, las cuatro semanas de cierre de la marcha blanca de la Etapa 2 coinciden con el peak de Fiestas Patrias; por eso toda habilitación de la Etapa 2 ocurre en agosto de 2028 y el cierre se mide con el volumen del peak (Formulario T-18, sección 3.1). LafroX adopta febrero de 2027 como supuesto de calendario de la carta Gantt.

Durante un mes de congelamiento que cae dentro de una marcha blanca no se inicia ninguna ola ni se despliega ningún cambio. La marcha blanca continúa en convivencia con la forma actual de trabajar, se siguen midiendo los indicadores diarios y el volumen real de septiembre sirve como prueba de carga en operación. La fecha efectiva de inicio se confirma con la consulta V-12 (Capítulo 3, Anexo 3.H).

\section{Anexo 7.B — Dependencias entre paquetes de trabajo}

La Tabla \ref{tab:7B1} lista las dependencias que estructuran la red del cronograma. Son de tipo fin-comienzo (FC), salvo las marcadas como comienzo-comienzo (CC) y las de tipo «por interfaz», en que el sucesor espera los contratos del predecesor para construir y su entrega para integrar; cada una indica su fundamento. El Formulario T-15, sección 6.1, aplica estas dependencias actividad por actividad. Sobre esta red se identifica la ruta crítica del Subdocumento 7, sección 7.3.1, y se aplica el método de programación del Formulario T-15.

\parte{tablas/tab_7B1}

Las dependencias D-29 a D-31 forman la ruta crítica de la Etapa 2. D-02, D-04 y D-14 a D-25 forman la cadena casi crítica de la Etapa 1. D-07 a D-13 convergen en el H3. D-32 a D-34 enlazan la certificación de la Etapa 2 con su marcha blanca, aceptación y operación.

\section{Anexo 7.C — Momento de los resultados de aceptación del caso}

El Caso 02, capítulo 18, exige indicar en qué momento del cronograma se alcanza cada resultado y cómo se mide. La Tabla \ref{tab:7C1} ubica en el cronograma las metas y métodos del Capítulo 3, Anexo 3.J.

\parte{tablas/tab_7C1}

Los resultados R18-01 a R18-10 y R18-12/R18-14/R18-15/R18-16 tienen verificaciones en la marcha blanca de la Etapa 1 (14 resultados); R18-11 y R18-13 corresponden a E2. Uno de ellos, R18-06, se prueba además antes del H5, y R18-12 se acepta en esa marcha blanca solo de forma provisional. Por eso la certificación de usuarios y el volumen real de esa marcha blanca concentran el riesgo de aceptación del proyecto. Dos resultados se confirman en operación: el tramo de OTIF del mes 32 y la pérdida de envases tras 12 meses.

\section{Anexo 7.D — Innovaciones en la EDT y en el cronograma}

Cada innovación declara sus paquetes de la EDT y sus meses del cronograma, como exige el Art. 29°, punto 4, de las Bases Administrativas. La Tabla \ref{tab:7D1} los presenta con su código de innovación y sus meses.

\parte{tablas/tab_7D1}

Las innovaciones 1 a 3 concentran su construcción entre los meses 2 y 10 y se validan en la marcha blanca de la Etapa 1, por lo que su riesgo queda acotado antes del H7. La innovación 5 sigue el calendario de la Etapa 2 y se extiende a las rutas durante los primeros meses de operación, y la innovación 4 liquida desde el mes 24, después de su línea base y de tres liquidaciones en sombra.

\section{Anexo 7.E — Trazabilidad de módulos e interfaces a la EDT}

La EDT contiene la arquitectura cuando cada módulo y cada interfaz del Capítulo 4 tienen al menos un paquete que los construye y prueba. La Tabla \ref{tab:7E1} presenta esa correspondencia para los doce módulos, y la Tabla \ref{tab:7E2}, para las quince interfaces.

\parte{tablas/tab_7E1}

Once de los doce módulos se construyen en la cuenta 3.4 y se prueban en la 3.8; M10 Analítica es el único módulo con trabajo en ambas etapas, porque sus indicadores operacionales entran en la Etapa 1 y el costo de servir en la Etapa 2.

\parte{tablas/tab_7E2}

Las quince interfaces tienen un paquete que las construye, y las interfaces internas (INT-01 a INT-05 e INT-12 a INT-14) dependen de la base compartida o de los módulos que las usan. Además de estas interfaces, el paquete 3.6.1 construye el intercambio de lote y trazabilidad con los proveedores.

\section{Anexo 7.F — Supuestos de planificación y dependencias para la gestión de riesgos}

Esta matriz traslada al Capítulo 8 los supuestos de la planificación que condicionan el cronograma, con su responsable y el control que los vigila. El Capítulo 8 los evalúa con su escala y los registra en el Formulario T-16. Los códigos P7 son identificadores de esta matriz y no renumeran el Formulario T-12.

\parte{tablas/tab_7F1}

Los once supuestos se concentran en tres frentes: el tamaño de las estimaciones (P7-01 a P7-03), la fecha efectiva y las ventanas de corte (P7-04 a P7-06) y las dependencias técnicas con la arquitectura (P7-08 a P7-10). Cada uno tiene un control que se revisa en el Comité de Proyecto, y el Capítulo 8 asigna a los de mayor exposición una reserva de contingencia.

\subsection{Política de reservas que recibe el Capítulo 8}

El modelo protege 3.072 HH para correcciones de la Etapa 1 en los meses 13 a 20, además del trabajo base y de la cobertura de servicio. El Formulario T-15, Tabla 5.2, fija reservas en días hábiles: H1 6, H2 35, H3 19, H4 19, H5 35, H8 13, H9 21 y H10 28. La subsanación de observaciones del CLIENTE consume la reserva del hito (T-15, sección 5.5). Toda amenaza se escala cuando la desviación proyectada supera la mitad de su reserva, sin esperar a consumirla. La holgura local de convergencia no se suma como reserva adicional. Las últimas cuatro semanas de marcha blanca no son reserva. La reserva de gestión se define en la Oferta Económica, porque la Oferta Técnica no contiene precios.

\subsection{Condiciones para aprobar la línea base}

Antes de aprobar la línea base del cronograma en el H1, el Jefe de Proyecto presenta al Comité de Proyecto la estimación validada por los líderes de frente, la red detallada, las personas y turnos asignados a cada frente, la fecha efectiva de inicio confirmada por el CLIENTE y las decisiones de la consulta V-12. La evidencia de los ensayos se incorpora a medida que cada hito la produce.

\section{Referencias}

Bases Administrativas, Art. 17°/18° y Formularios T-14/T-15/T-18; Bases Técnicas Transversales, RT-20.02/05/06; Caso 02, §13.3/§17.6; Aclaraciones, Capítulos 7/8; SD3 §3.4.4; SD4 y anexos lógicos; SD6 §6.1/§6.2; formularios del SD7.

\section{Declaración de uso de IA}

En cumplimiento de la sección 7.2 de las Aclaraciones de la licitación, la tabla siguiente declara el uso de herramientas de inteligencia artificial en estos anexos, con la revisión humana de cada parte. La declaración se consolida en el Formulario A-6.

\parte{tablas/tab_IA}

